\section{Introduction} \label{sec:intro}
Understanding the mechanisms underlying moral reasoning is central to a broad range of
research endeavors in Natural Language Processing. Morality has been recognized as a
source of disagreement among humans (cit, cit) and its automatic detection (cit, cit)
supports a better understanding of how political and social views are shaped in the
public debate (cit, cit). Moral reasoning is also a key component of research on
Human-AI value alignment (cit, cit), which examines whether and how Large Language
Models (LLMs) represent, and can be steered toward, different moral perspectives
(cit, cit).

NLP research on moral reasoning relies on two psychological theories of value and
moral pluralism (cit), which provide established taxonomies to elicit conflicting
moral views: Schwartz's theory of basic human values (cit) and the Moral Foundations
Theory (MFT) (cit). A systematic operationalization of these two theories is still
missing, and this results in scattered research outputs and directions. A
complementary problem is methodological: adopting tests from social science to
benchmark NLP technologies has been shown to lack reliability (cit) and
reproducibility (cit). Questionnaire scores administered to LLMs are sensitive to
prompt phrasing and option order, and a score alone offers no evidence that the model
is applying the construct the instrument was designed to measure. Any framework for
measuring morality in LLMs must therefore treat the validity of its own instruments as
part of what is being tested, rather than as an assumption.

Our work addresses this challenge by providing the first comprehensive evaluation of
LLM moral profiles under the lens of both Schwartz's values and MFT. Grounding our
study in psychological research, we structure the evaluation in two stages, each
designed to test a distinct aspect of measurement validity. The \textbf{internal
validation} asks whether questionnaire responses mean what they appear to mean: we
combine quantitative and qualitative methods to assess the coherence between how
models rank moral claims and how they justify those same claims in open-ended
interviews, and we characterise the resulting differences in moral profiles.
Divergence between ranking and justification is itself a diagnostic, indicating where
scores should not be read at face value. The \textbf{external validation} asks whether
the profiles have consequences beyond the instrument that produced them: we steer
models with moral opinions elicited through interviews and measure their effect on
Hate Speech detection. A profile that predicts downstream behaviour is measuring
something; one that does not is an artefact of the questionnaire.

Our results show that, this, and that.

Our contribution is the following:
\begin{itemize}
    \item we present a framework grounded in social psychology for the study of LLM
    moral profiles, in which the reliability of the psychological instruments is
    evaluated alongside the models themselves;
    \item we release a corpus of LLM interviews based on psychological questionnaires,
    human-annotated for [X], and aligned with the corresponding questionnaire rankings;
    \item we systematically measure the steering effect of psychologically-validated
    items in a real-world scenario and its differential impact across models.
\end{itemize}
%Schwartz wheel of values (cit), which organizes 19 values as adjacent (e.g., hedonism and achievement) or opposite (e.g., hedonism and benevolence); the Moral Foundations Theory


\section{Experimental Setup} \label{sec:setup}

\subsection{Models} \label{sec:setup-models}
We evaluate six open-weight instruction-tuned LLMs in the 7--8B parameter range:
Apertus-8B-Instruct, Falcon3-7B-Instruct, Llama-3.1-8B-Instruct,
Ministral-3-8B-Instruct(-2512), Olmo-3-7B-Instruct, and Qwen3-8B (cit, cit, cit, cit,
cit, cit). We restrict our study to this size class so that any differences we observe
reflect training data, alignment procedure, and architecture rather than scale; we
leave a scale-controlled extension to future work (\S\ref{sec:discussion}).

\subsection{Instruments} \label{sec:setup-instruments}
We elicit moral profiles with two validated psychometric instruments. The
\textbf{Moral Foundations Questionnaire} (MFT; cit) contains 36 items spanning six
foundations (\emph{care}, \emph{equality}, \emph{proportionality}, \emph{loyalty},
\emph{authority}, \emph{purity}), each rated on a 5-point Likert scale. The
\textbf{Portrait Values Questionnaire} (PVQ; cit) contains 40 items spanning
Schwartz's ten basic values (\emph{universalism}, \emph{benevolence}, \emph{tradition},
\emph{conformity}, \emph{security}, \emph{power}, \emph{achievement}, \emph{hedonism},
\emph{stimulation}, \emph{self-direction}), each rated on a 6-point Likert scale. Both
instruments present short first- or third-person statements (e.g., ``I think people
should be rewarded in proportion to what they contribute'') that respondents rate for
agreement or self-similarity.

\subsection{Internal Validation Protocol} \label{sec:setup-internal}
For every item of both instruments, we elicit two responses from each model under an
identical framing (participant in a psychological survey): (i) a \textbf{Likert
rating} on the instrument's native scale (\texttt{opinion}), obtained with constrained
decoding over the label set; and (ii) a \textbf{free-text justification}
(\texttt{belief}) obtained by prompting the model to explain its position on the same
statement in open-ended language, with no numeric constraint. The Likert rating
operationalises how a model \emph{ranks} a moral claim; the free-text justification
operationalises how it \emph{justifies} that same claim, and is the artefact we later
present to human raters (internal validation) and use to steer a downstream classifier
(external validation, \S\ref{sec:setup-external}).

To assess whether the free-text justification is consistent with the Likert rating, we
recruit $N=54$ crowd-annotators (Prolific), stratified $3\times3\times3$ within every
model by gender (woman / man / non-binary) and continent of birth (Africa / Asia /
Europe), giving nine annotators per model. Each annotator reads a model's free-text
justification for every item of one instrument and rates their agreement with it on
the instrument's native scale, blind to the model's own Likert rating. For item $i$ and
model $m$ we compute the \textbf{external opinion} $\bar{y}_{i,m}$ as the mean of the
nine annotator ratings, and the \textbf{annotator disagreement} $d_{i,m}$ as the mean
absolute deviation of those ratings around their own mean. Internal validity is
operationalised as the correlation between a model's own Likert rating and its
external opinion: high, well-calibrated agreement indicates the free-text
justification says what the Likert score implies it should say; low agreement
indicates the two elicitation formats are measuring different things, or that the
Likert score does not reflect the model's stated reasoning.

We separately validate the annotation protocol itself, following the principle
(\S\ref{sec:intro}) that instrument reliability must be evaluated alongside the
models: we test (a) whether an annotator's own reply to an item predicts their
evaluation of a model's reply to the same item (a possible source of rater bias), and
(b) whether annotators sharing a demographic attribute (gender, continent of birth)
agree with each other more than annotators who do not.

\subsection{External Validation Protocol} \label{sec:setup-external}
We test whether the free-text justifications elicited in \S\ref{sec:setup-internal}
have measurable consequences beyond the instrument that produced them, using implicit
hate speech detection as the downstream task (cit). We sample $N=500$ messages from the
Implicit Hate Corpus (cit), balanced 250 hate / 250 non-hate, and prompt each of the
six models to classify each message as a strict binary decision under three
conditions: (i) \textbf{zero-shot}, with no additional context; and (ii) \textbf{belief-steered},
once per questionnaire item (76 runs per model: 36 MFT + 40 PVQ), in which the
model's own free-text justification for that item is prepended to the classification
prompt as context (``Beliefs: \{belief\}''). Decoding is greedy (temperature 0) and
constrained to the label set $\{0,1\}$. Critically, steering is \emph{paired by model}:
a model is only ever steered by its own justifications, never by another model's, so
any difference we observe between models isolates how differently each model responds
to being shown its own stated beliefs, not a difference in whose beliefs are shown.
Our primary metric is \textbf{recall} on the hate-speech class (proportion of the 250
positive messages correctly flagged), since a shift in recall is the safety-relevant
direction of error for a content-moderation classifier.

\subsection{Statistical Analysis} \label{sec:setup-stats}
Because Likert ratings and recall are ordinal/discrete outcomes on a scale that is
externally fixed and shared across all models (rather than a subjective, rater-specific
anchor), we treat a model's absolute response level as a substantive quantity rather
than nuisance variance, and we do not normalise ratings across models before testing
between-model differences. We use Pearson and Spearman correlation to test the
strength of the association between two continuous/ordinal measures; a two-sample
Mann-Whitney $U$ test and its $k$-sample generalisation, the Kruskal-Wallis $H$ test,
to compare the distribution of a measure between models (both invariant to how the
scale is used and preferable to a $\chi^2$ test of independence, which treats ordinal
categories as unordered nominal ones); and an exact (binomial) McNemar test for the
paired binary outcome of message-level classification correctness under zero-shot
versus steered conditions. Where we report many tests within one hypothesis family
(e.g., all items within one model), we control the false discovery rate with the
Benjamini-Hochberg procedure (cit) and report both raw and corrected $p$-values.


\section{Results} \label{sec:results}

\subsection{Internal Validation} \label{sec:results-internal}

\paragraph{Ranking and justification are only weakly coherent.}
Correlating each model's Likert rating against its external opinion
(\S\ref{sec:setup-internal}) yields a significant positive correlation for only
2 of 6 models on MFT (Apertus, $r=0.60$; Ministral, $r=0.39$) and 1 of 6 on PVQ
(Qwen, $r=0.58$); the remaining models show no reliable relationship between how they
rank an item and how a human rates the justification they give for that ranking. This
is the central internal-validity finding of our study: a Likert score is, for most
models, not a reliable summary of the model's own stated reasoning.

\paragraph{Models differ far more in absolute response level than in relative
ranking.} A Kruskal-Wallis test on raw Likert ratings across the six models is highly
significant on both instruments, with a large effect size ($\eta^2 \approx 0.45$--$0.46$
on MFT and PVQ respectively) — model identity accounts for roughly half of the
variance in ratings. One model (Llama, on PVQ) returned the identical rating on all 40
items, exhibiting no item-level discrimination at all. External opinion (the human
rating of the justification) also differs significantly by model, but with a smaller
effect ($\eta^2 \approx 0.11$--$0.14$).

\paragraph{Human ratings of justifications track item content, not model identity.}
Despite the level differences above, the \emph{pattern} of which items draw a higher
or lower external opinion is highly consistent across the six models' justifications:
pairwise correlation of external opinion between every pair of models has a median
$r=0.76$ on PVQ (all 15 pairs significant, $p<10^{-5}$) and a weaker but uniformly
positive median $r=0.29$ on MFT. A distribution-free test of whether this pattern's
\emph{shape} depends on model identity is non-significant on both instruments. Taken
together, this indicates that what varies by model is the average level of agreement
elicited by the free-text justifications, not which items are relatively more or less
convincing — that ordering is largely a property of the item itself.

\paragraph{The annotation instrument is reliable.} An annotator's own reply to an item
correlates only modestly with their evaluation of a model's reply to that same item
(pooled $r=0.31$ on MFT, $r=0.40$ on PVQ; $R^2 \approx 10$--$16\%$), indicating that
personal stance is a minor, not a dominant, contributor to how an annotator rates a
justification. We find no reliable effect of annotator demographics on inter-annotator
agreement: gender shows no effect on either instrument (all $p \geq 0.74$, properly
powered at 54 same-gender vs. 162 different-gender pairs), and continent of birth
reaches $p=0.031$ on one of four MFT metrics but does not replicate on PVQ ($p=0.50$),
consistent with chance given the eight pooled tests run.

\subsection{External Validation} \label{sec:results-external}

\paragraph{Steering with a model's own beliefs shifts recall, and the direction is
model-specific.} A one-sample Wilcoxon signed-rank test on the 76 per-item recall
deltas (steered $-$ zero-shot) is significant for 5 of 6 models
(Table~\ref{tab:steering-effect}). Steering \emph{increases} recall for Ministral
($+0.090$), Falcon ($+0.071$), Qwen ($+0.051$), and Llama ($+0.040$); it
\emph{decreases} recall for Olmo ($-0.096$), of comparable magnitude to Ministral's
gain; Apertus shows no net average effect. A Kruskal-Wallis test on the per-item delta
confirms the difference between models is itself large ($\eta^2 \approx 0.5$), and
every one of the 15 pairwise Mann-Whitney comparisons between models is significant.

\begin{table}[t]
\centering
\caption{Recall under zero-shot and belief-steered conditions, per model (mean over 76
items). $\ast$ marks a significant Wilcoxon signed-rank test on the per-item delta
($p<0.05$).}
\label{tab:steering-effect}
\begin{tabular}{lccc}
\toprule
Model & Zero-shot recall & Steered recall (mean) & $\Delta$ \\
\midrule
Ministral-3-8B-Instruct  & 0.856 & 0.946 & $+0.090^{\ast}$ \\
Falcon3-7B-Instruct      & 0.620 & 0.691 & $+0.071^{\ast}$ \\
Qwen3-8B                 & 0.716 & 0.767 & $+0.051^{\ast}$ \\
Llama-3.1-8B-Instruct    & 0.896 & 0.936 & $+0.040^{\ast}$ \\
Apertus-8B-Instruct      & 0.856 & 0.853 & $-0.003$ \\
Olmo-3-7B-Instruct       & 0.628 & 0.532 & $-0.096^{\ast}$ \\
\bottomrule
\end{tabular}
\end{table}

\paragraph{Item-level effects are individually significant, and reveal that
``no net effect'' is not the same as ``unaffected''.} The aggregate test above has only
$n=6$ models per item and is underpowered at the item level (0 of 76 items survive
Benjamini-Hochberg correction). We therefore repeat the analysis at the message level
within each model: for each model and item, an exact McNemar test on the 250
positive-class messages (paired zero-shot vs.\ steered correctness), corrected within
each model's own 76-item family, has far more power. Between 23 (Apertus) and 74
(Ministral) of the 76 items individually reach significance per model. For Falcon,
Llama, Ministral, and Qwen, essentially every significant item is an \emph{increase}
in recall (53/56, 51/51, 74/74, and 45/45 respectively); for Olmo, most significant
items (35/46) are \emph{decreases}. Apertus, in contrast to its near-zero average, is
genuinely mixed at the item level (13 decreases, 10 increases): its average was masking
real, opposite-signed effects rather than reflecting an absence of effect.

\paragraph{The content that drives the effect replicates across two independent
instruments.} Aggregating item-level results by content category, rather than testing
single items, restores statistical power and reveals a convergent pattern
(Table~\ref{tab:category-patterns}). On PVQ, every Schwartz value has a net positive
pooled effect on recall \emph{except} universalism ($-0.048$) and benevolence
($-0.044$) — the two values in Schwartz's self-transcendence quadrant, both concerned
with the welfare of others. On MFT, \emph{care} ($+0.017$) and \emph{purity}
($+0.007$) are the two weakest foundations by a wide margin, well behind
proportionality ($+0.052$), loyalty ($+0.050$), equality ($+0.048$), and authority
($+0.034$); \emph{care} is MFT's closest conceptual analogue to PVQ's
universalism/benevolence. In both instruments, independently, the same two models —
Apertus and Olmo — are the only ones for which this content category has a
\emph{negative} mean delta (Apertus: universalism $-0.032$, benevolence $-0.061$, care
$-0.041$, purity $-0.031$; Olmo: universalism $-0.391$, benevolence $-0.320$, care
$-0.022$, purity $-0.069$); the remaining four models are helped by this content too,
only less than by any other category.

\begin{table}[t]
\centering
\caption{Mean recall delta (steered $-$ zero-shot) pooled across all six models, by
content category. $n_{\text{sig}}$: number of (model, item) pairs individually
significant after within-model FDR correction, out of the category's item count
$\times$ 6 models.}
\label{tab:category-patterns}
\begin{tabular}{lrr}
\toprule
Category & Mean $\Delta$ recall & $n_{\text{sig}}$ \\
\midrule
\multicolumn{3}{l}{\emph{PVQ (Schwartz values)}} \\
Power           & $+0.059$ & 9/18 \\
Stimulation     & $+0.057$ & 14/18 \\
Security        & $+0.052$ & 22/30 \\
Achievement     & $+0.050$ & 16/24 \\
Self-Direction  & $+0.036$ & 15/24 \\
Tradition       & $+0.026$ & 16/24 \\
Hedonism        & $+0.020$ & 14/18 \\
Conformity      & $+0.015$ & 13/24 \\
Benevolence     & $-0.044$ & 15/24 \\
Universalism    & $-0.048$ & 24/36 \\
\midrule
\multicolumn{3}{l}{\emph{MFT (Moral Foundations)}} \\
Proportionality & $+0.052$ & 28/36 \\
Loyalty         & $+0.050$ & 27/36 \\
Equality        & $+0.048$ & 31/36 \\
Authority       & $+0.034$ & 25/36 \\
Care            & $+0.017$ & 11/36 \\
Purity          & $+0.007$ & 15/36 \\
\bottomrule
\end{tabular}
\end{table}


\section{Discussion} \label{sec:discussion}

\paragraph{Likert scores alone are not a trustworthy moral profile.} Our internal
validation shows that, for four of six models, a model's Likert rating and the human
reading of the reasoning it gives for that rating are statistically independent. This
substantiates the methodological concern raised in \S\ref{sec:intro}: reporting a
Likert-scale moral profile without checking it against the model's own justification
risks describing an artefact of the instrument rather than a property of the model.
The large between-model effect on absolute rating level, contrasted with the
comparatively small, item-driven variation in how justifications are received,
suggests that much of what a Likert-only study would report as ``moral difference
between models'' is actually a difference in how liberally each model uses the rating
scale — a calibration effect, not necessarily a difference in the underlying
construct.

\paragraph{Free-text justifications are behaviourally real, even where Likert scores
are not.} The external validation result is, in this light, the more informative one:
despite the weak internal coherence documented above, the free-text justifications
have a large, consistent, and — at the category level — cross-instrument-replicated
effect on a completely unrelated downstream task. This is precisely the diagnostic
proposed in \S\ref{sec:intro}: a profile that predicts downstream behaviour is
measuring something, even when it fails an internal consistency check against its own
Likert summary. We read this as evidence that the free-text elicitation captures a
real behavioural disposition in the model — one that a Likert score, taken alone, would
have missed or mischaracterised.

\paragraph{A convergent, model-specific vulnerability.} The clearest actionable
finding is that Apertus and Olmo are, independently on two conceptually distinct
instruments (PVQ's universalism/benevolence and MFT's care/purity, both centred on
concern for others' welfare), the only models whose hate-speech recall is
\emph{hurt} by being shown their own other-regarding moral reasoning, while the same
content helps the remaining four models, if less than other content does. We do not
have access to these models' training data or alignment procedures and so cannot
adjudicate between candidate mechanisms — e.g., a shift in decision threshold induced
by sympathetic-sounding context, versus a more structural interaction between
other-regarding framing and the specific classifier prompt — and we flag this
explicitly as a hypothesis for follow-up mechanistic work rather than an established
explanation.

\paragraph{Limitations.} Our study covers six models in a single parameter range, one
downstream task, and one hate-speech corpus; we do not yet have a neutral-content
control condition to isolate the effect of moral \emph{content} from the effect of
prepending \emph{any} plausible-sounding context, nor repeated sampling at
non-zero temperature to estimate generation variance. The item-level cross-model test
(\S\ref{sec:results-external}) is underpowered by construction ($n=6$); we mitigate
this by aggregating to the content-category level, but single-item claims should be
read as descriptive, not confirmatory. Finally, our steering condition is paired by
model (a model sees only its own justifications): this isolates how differently
models respond to their own stated beliefs, but does not test whether the effect
transfers when a model is steered by another model's beliefs, which we leave to future
work.

\paragraph{Future work.} A natural extension is to scale the model pool (in count and
in parameter size) to test whether the internal-validity gap and the
universalism/benevolence vulnerability are properties of this size class or of
particular training regimes; to add a neutral-content control and a second downstream
task/corpus to isolate content-specific effects from generic context effects; and to
probe the mechanism behind the Apertus/Olmo result directly, for instance via
decision-threshold or logit-level analysis rather than hard-label outcomes alone.
